% from main.tex
%%%%%%%%%%  scan idx 277  |  folio 264  %%%%%%%%%%
Pour tout changement de base \(Y'\to Y\), la flèche canonique
\[
 R^pf_*(E)\otimes_Y\mathcal O_{Y'}\longrightarrow H^p(Rf_*(E)\otimes_Y^{\mathbf L}\mathcal O_{Y'})
\]
est un isomorphisme.
\end{enumerate}
Alors on a :
\[
 \mathrm a)\Longleftrightarrow\mathrm a')\Longleftrightarrow\mathrm b)\Longleftrightarrow\mathrm c)\Longleftarrow\mathrm b').
\]
En outre, si \(f\) est propre, \(E\) pseudo-cohérent rel. à \(f\) (donc parfait rel. à \(f\)), et sous réserve de validité de (2.1) pour \(f\) (donc en particulier si \(f\) est projectif ou si \(Y\) est noethérien), les conditions précédentes sont équivalentes.

\emph{Preuve.} Les implications \(\mathrm a)\Longleftrightarrow\mathrm a')\Longleftrightarrow\mathrm b)\Longleftarrow\mathrm b')\) résultent aussitôt de (3.7), (5.1) et (5.2 (iv)), de même que l'implication \(\mathrm b')\Longleftarrow\mathrm b)\) (pour \(f\) propre et \(E\) pseudo-cohérent rel. à \(f\)) en vertu de (2.1). Il reste à prouver qu'on a \(\mathrm a)\Longleftrightarrow\mathrm c)\). Cela va découler de la première assertion du lemme suivant :

\textbf{Lemme 5.3.1.} Soit \(X\) un schéma quasi-séparé, et soit \(E\in\Ob D^{-}(X)\) un complexe à cohomologie localement bornée et quasi-cohérente. Alors la condition (5.1 a)) équivaut à la condition
\begin{enumerate}[label=\alph*)]
\item[c)] Pour tout schéma \(X'\) au-dessus de \(X\) la flèche canonique
\[
 H^p(E)\otimes_X\mathcal O_{X'}\longrightarrow H^p(E\otimes_X^{\mathbf L}\mathcal O_{X'})
\]
est un isomorphisme.
\end{enumerate}
En outre, si \(E\) est pseudo-cohérent, la condition (5.1 a)) implique que, pour tout
